% !TeX root = ../main.tex
% Section 2: The Weyl-Heisenberg Algebra of Market Observables
% v4.1 (2026-08-30): Added §2.0 "Microstructural Origin of the CCR" before §2.1.
% Directly addresses MF handling-editor request: "it is not clear why eq (5)
% is a good model for financial prices." Abstract axis-frozen; §2 structure
% preserved; only new subsection prepended.

\section{The Weyl--Heisenberg Algebra of Market Observables}\label{sec:weyl-heisenberg}

% -----------------------------------------------------------------------
\subsection{Microstructural Origin of the Canonical Commutation Relation}
\label{sec:ccr-derivation}
% -----------------------------------------------------------------------

Equation~\eqref{eq:ccr-intro} is not imposed as a modeling assumption but
\emph{derived} from Kyle's (1985) discrete-time price-impact model.
Consider a single-period Kyle market with a risk-neutral market maker, an
informed trader submitting order $q_t$, and noise traders. The market maker
sets price as
\begin{equation}\label{eq:kyle-price-impact}
    X_{t+\Delta t} = X_t + \lambda \, q_t + \varepsilon_t,
\end{equation}
where $\lambda > 0$ is the price-impact coefficient and $\varepsilon_t$ is
exogenous noise \citep{kyle1985continuous}. Define two operators on the
trader's position state space $L^2(\mathbb{R})$:
\begin{equation}\label{eq:ops-def}
    \hat{X}_{\mathrm{obs}} f(q) := (X_t + \lambda q)\, f(q), \qquad
    \hat{P}_{\mathrm{net}} f(q) := -i\sigma_\varepsilon\sqrt{\Delta t}
    \,\frac{\partial f}{\partial q}(q),
\end{equation}
where $f \in L^2(\mathbb{R})$ is the trader's position wavefunction.
$\hat{X}_{\mathrm{obs}}$ represents the \emph{post-trade} observable price
$X_t + \lambda q$ (which depends on the order $q$ through the price-impact
rule~\eqref{eq:kyle-price-impact}); $\hat{P}_{\mathrm{net}}$ is the
infinitesimal order-flow generator. The \emph{observe-then-trade} sequence
first applies $\hat{P}_{\mathrm{net}}$ (shifting position by a unit order)
then reads $\hat{X}_{\mathrm{obs}}$ (recording the resulting price); the
\emph{trade-then-observe} sequence applies them in reverse order.
The two operator compositions are
\begin{align}
    (\hat{X}_{\mathrm{obs}}\,\hat{P}_{\mathrm{net}}) f(q)
    &= (X_t + \lambda q)\!\left(-i\sigma_\varepsilon\sqrt{\Delta t}\,f'(q)\right), \notag\\
    (\hat{P}_{\mathrm{net}}\,\hat{X}_{\mathrm{obs}}) f(q)
    &= -i\sigma_\varepsilon\sqrt{\Delta t}\,\frac{\partial}{\partial q}
       \!\bigl[(X_t + \lambda q)\,f(q)\bigr]
    = -i\sigma_\varepsilon\sqrt{\Delta t}\bigl[(X_t + \lambda q)\,f'(q)
       + \lambda f(q)\bigr],
    \label{eq:composition}
\end{align}
where the second line uses the product rule and $\partial(\lambda q)/\partial q
= \lambda$, reflecting the linear price-impact rule~\eqref{eq:kyle-price-impact}.
Taking the difference yields
\begin{equation}\label{eq:ccr-derivation}
    [\hat{X}_{\mathrm{obs}},\,\hat{P}_{\mathrm{net}}]\,f(q)
    = \bigl(\hat{X}_{\mathrm{obs}}\hat{P}_{\mathrm{net}}
      - \hat{P}_{\mathrm{net}}\hat{X}_{\mathrm{obs}}\bigr)f(q)
    = i\lambda\sigma_\varepsilon\sqrt{\Delta t}\,f(q),
\end{equation}
so that $[\hat{X}_{\mathrm{obs}},\hat{P}_{\mathrm{net}}]
= i\lambda\sigma_\varepsilon\sqrt{\Delta t}\,\mathbf{1}
=: i\hbar_{\mathrm{econ}}\,\mathbf{1}$,
with $\hbar_{\mathrm{econ}} := \lambda \cdot q_0$ and
$q_0 := \sigma_\varepsilon\sqrt{\Delta t}$ a reference order-flow unit.
This identifies $\hbar_{\mathrm{econ}}$ as directly estimable from the slope
of the empirical price-impact function
\citep{almgren2005direct,toth2011anomalous}.
In Kyle's model, $\lambda$ measures the permanent price impact per unit of
order flow. Economically, $\hbar_{\mathrm{econ}}$ quantifies the minimum
distinguishable increment in price--order-flow space: below this scale,
simultaneous observation of price and order flow is precluded by the trading
mechanism itself, paralleling Heisenberg's measurement bound in physical
systems. This connects the quantum friction scale directly to observable
market microstructure depth.

The continuous-time limit $\Delta t \to 0$ with $\lambda \sim \Delta t^{-1/2}$
(consistent with the Almgren--Chriss scaling \citep{almgren2005direct}) yields the
CCR \eqref{eq:ccr-intro}. Formally, as $\Delta t \to 0$ with
$\hbar_{\mathrm{econ}} = \lambda \sigma_\varepsilon \sqrt{\Delta t}$ held fixed,
the discrete operators $\hat{X}_{\mathrm{obs}}^{(\Delta t)}$ and
$\hat{P}_{\mathrm{net}}^{(\Delta t)}$ converge in the strong resolvent sense
to the canonical position and momentum operators on $L^2(\mathbb{R})$ satisfying
$[\hat{X}, \hat{P}] = i\hbar_{\mathrm{econ}}\mathbf{1}$ \citep{reed1972methods}.
The Trotter--Kato theorem \citep{kato1978trotter} ensures that the semigroup
generated by the limiting Lindblad equation is the \emph{strong} operator limit
(in the sense of pointwise convergence on each vector) of the
discrete-time Markov semigroups, providing a rigorous foundation for the
continuous-time framework. The frictionless limit $\lambda \to 0$ (or equivalently
$\hbar_{\mathrm{econ}} \to 0$) recovers a commutative algebra
$[\hat{X}, \hat{P}] = 0$, which is precisely the classical BSM assumption
\citep{black1973pricing,merton1973theory}. The CCR is therefore not a formal
analogy from physics but a \emph{derived consequence} of finite price impact in the
sense of Kyle \citep{kyle1985continuous} and Glosten--Milgrom \citep{glosten1985bid}.

\begin{proposition}[Non-commutativity is necessary under price impact]
\label{prop:classical-failure}
Suppose the price and order-flow operators $\hat{X}$ and $\hat{P}$ satisfy
Kyle's linear price-impact rule \eqref{eq:kyle-price-impact} with $\lambda > 0$.
Then:
\begin{enumerate}
    \item[(i)] $[\hat{X}, \hat{P}] = i\hbar_{\mathrm{econ}}\mathbf{1} \neq 0$
        (non-commutativity is forced by the market mechanism).
    \item[(ii)] If one instead imposes $[\hat{X}, \hat{P}] = 0$ (commutative
        algebra), then the operator $\hat{P}$ cannot generate bounded translations
        of $\hat{X}$ in $L^2(\mathbb{R})$, and no self-financing strategy
        $\delta_t \in \mathcal{K}$ satisfies the Kyle pricing rule for all order
        sizes $q \in \mathbb{R}$. Equivalently, perfect dynamic replication
        is unattainable for any payoff with non-trivial order-flow sensitivity.
\end{enumerate}
\end{proposition}

\begin{proof}
Part (i) is established in Section~\ref{sec:ccr-derivation}.
For part (ii): if $[\hat{X}, \hat{P}] = 0$, the Stone--von Neumann theorem
(Proposition~\ref{prop:stone-vn}) does not apply and the algebra $\mathcal{A}$
is commutative, i.e., isomorphic to $C_0(\mathbb{R}^2)$ (continuous functions
vanishing at infinity on phase space). In a commutative algebra, every element
is a function of commuting variables, so $\hat{P}$ cannot generate non-trivial
phase-space rotations. Concretely, the composition identity
$\hat{X}_{\mathrm{obs}}\hat{P}_{\mathrm{net}}f = \hat{P}_{\mathrm{net}}\hat{X}_{\mathrm{obs}}f$
would require $\lambda = 0$ (from \eqref{eq:ccr-derivation}), contradicting
$\lambda > 0$. Hence, under Kyle price impact, the commutative framework is
internally inconsistent.
\end{proof}

\begin{remark}[Why not a stochastic control formulation?]
\label{rem:why-not-sde}
An alternative response to price impact is to augment the classical It\^o SDE
with a feedback term $dX_t = \sigma dW_t + \lambda \delta_t\,dt$, where
$\delta_t$ is the hedge ratio. This approach---used in Almgren-Chriss (2001)
and related execution-cost models---treats price impact as an exogenous cost
function and preserves the commutative probability space. The key distinction
is that such models assume the hedge ratio $\delta_t$ and the post-trade
price $X_t + \lambda\delta_t$ can be simultaneously observed: $\delta_t$
and $X_t$ commute as random variables. Proposition~\ref{prop:classical-failure}
shows this simultaneous observability fails when price impact is linear in
order flow: observing the post-trade price already incorporates the order,
and the order in turn responds to the pre-trade price, generating irreducible
non-commutativity at the level of the algebra of observables rather than at the
level of stochastic paths.
\end{remark}

\begin{remark}[Continuous-time scaling and Back's equilibrium]
\label{rem:back-scaling}
The scaling $\lambda \sim \Delta t^{-1/2}$ used above follows the Almgren--Chriss
execution-cost framework \citep{almgren2005direct}, in which market impact is an
\emph{exogenous} cost function. In the continuous-time generalisation of Kyle's
equilibrium \citep{kyle1985continuous}, the endogenous price-impact coefficient satisfies
$\lambda_t = \sigma_v / (2\sigma_u \sqrt{T-t})$, which has a different time-dependence.
The present paper uses the Almgren--Chriss scaling to fix $\hbar_{\mathrm{econ}}$
as a time-invariant parameter; this is a modelling choice that simplifies the empirical
calibration at the cost of abstracting from the time-decay of $\lambda_t$ in Back's
equilibrium. The calibrated value $\hbar_{\mathrm{econ}} = 0.087$ should therefore
be interpreted as a \emph{time-averaged} friction parameter over the calibration
window, rather than an instantaneous equilibrium coefficient.
\end{remark}

\begin{remark}[Interpretation]
Equation~\eqref{eq:ccr-derivation} shows that the post-trade price and the
magnitude of the executed order (operators defined in~\eqref{eq:ops-def}) cannot
be simultaneously resolved with arbitrary
precision in any market with $\lambda > 0$. This is the financial analogue of
the Heisenberg uncertainty principle: measuring $X_t$ precisely requires a
transaction that perturbs $P_t$, and vice versa. The economic Planck constant
$\hbar_{\mathrm{econ}}$ quantifies the irreducible cost of microstructural
friction, and vanishes only in the frictionless limit $\lambda \to 0$.
\end{remark}


\begin{definition}[Economic Planck constant]\label{def:hbar-econ}
Let $\hat{X}$ be the log-price operator and $\hat{P}$ be the order-flow momentum
operator, defined such that $\hat{P}$ generates translations in $\hat{X}$ via the
Weyl commutation relation \citep{Weyl1927}. The \emph{economic Planck constant}
is defined as
\begin{equation}\label{eq:hbar-def}
    \hbar_{\mathrm{econ}} := \lambda \cdot q_0
    = \lambda \cdot \sigma_\varepsilon \sqrt{\Delta t},
\end{equation}
where $\lambda > 0$ is Kyle's price-impact coefficient, $\sigma_\varepsilon$ is
the standard deviation of order-flow noise, and $\Delta t$ is the minimum
trading time scale. By construction, $\hbar_{\mathrm{econ}}$ satisfies
\begin{equation}\label{eq:ccr}
    [\hat{X}, \hat{P}] = i\hbar_{\mathrm{econ}} \mathbf{1},
\end{equation}
where $\mathbf{1}$ is the identity operator on the market Hilbert space, as
derived in Section~\ref{sec:ccr-derivation}.
\end{definition}

The commutation relation \eqref{eq:ccr} is the mathematical expression of the fact
that price and order flow are complementary observables in a market with microstructure
frictions; $\hbar_{\mathrm{econ}}$ can be estimated directly from the slope of the
empirical price impact function (Remark~\ref{rem:units}).

\begin{remark}[Units and interpretation]\label{rem:units}
In dimensional terms, $\hbar_{\mathrm{econ}} = \lambda \cdot q_0$ carries the
units of price $\times$ order-flow, i.e., $[\text{price}] \times [\text{shares}]$.
In the log-return / normalized-order-flow parametrization used throughout this
paper---where prices are expressed as log-returns and order flow is normalized by
$\sigma_\varepsilon\sqrt{\Delta t}$---both observables become dimensionless,
so $\hbar_{\mathrm{econ}}$ is effectively dimensionless in this convention.
Typical values estimated from market data range from $0.01$ to $0.15$,
with higher values corresponding to more illiquid markets.
\end{remark}

\begin{remark}[Notation]\label{rem:notation}
We use $\hbar_{\mathrm{econ}}$ as the canonical symbol throughout.
\end{remark}

\begin{theorem}[Financial uncertainty principle]\label{thm:uncertainty}
Let $\hat{X}$ be the log-price operator and $\hat{P}$ be the order-flow momentum
operator, satisfying the canonical commutation relation
$[\hat{X}, \hat{P}] = i\hbar_{\mathrm{econ}} \mathbf{1}$. Then for any market state
$\omega$, the variances of $\hat{X}$ and $\hat{P}$ satisfy
\begin{equation*}
    \sigma_X \cdot \sigma_P \geq \frac{\hbar_{\mathrm{econ}}}{2},
\end{equation*}
where $\sigma_X^2 = \omega(\hat{X}^2) - \omega(\hat{X})^2$ and
$\sigma_P^2 = \omega(\hat{P}^2) - \omega(\hat{P})^2$.
\end{theorem}

\begin{proof}
Since $\hat{X}$ and $\hat{P}$ are unbounded operators, the Robertson
inequality applies on the dense domain
$\mathcal{D} := \mathcal{D}(\hat{X}\hat{P}) \cap \mathcal{D}(\hat{P}\hat{X})
\subset \mathcal{H}_\omega$, which is non-empty and dense in the Schr\"{o}dinger
representation by the Stone--von Neumann theorem \citep{vonNeumann1931}. For
any $|\psi\rangle \in \mathcal{D}$, the Cauchy--Schwarz inequality applied
to the commutator $[\hat{X}, \hat{P}]$ in the GNS representation gives
\begin{equation*}
    \|(\hat{X} - \langle \hat{X} \rangle)|\psi\rangle\| \cdot
    \|(\hat{P} - \langle \hat{P} \rangle)|\psi\rangle\|
    \geq \frac{1}{2} |\langle \psi | [\hat{X}, \hat{P}] | \psi \rangle|
    = \frac{\hbar_{\mathrm{econ}}}{2}.
\end{equation*}
Taking the infimum over all $|\psi\rangle \in \mathcal{D}$ with
$\||\psi\rangle\| = 1$ and identifying $\sigma_X^2, \sigma_P^2$ with the
respective variances in state $\omega$ yields the inequality for all states
$\omega$ whose GNS vector lies in $\mathcal{D}$.
\end{proof}

\subsection{The Weyl Form of the Canonical Commutation Relation}

Equation \eqref{eq:ccr} is unbounded and presents domain issues in infinite
dimensions. The Weyl form exponentiates the unbounded operators to unitary operators,
yielding a rigorous algebraic formulation.

\begin{definition}[Weyl operators]\label{def:weyl-ops}
For each $(a, b) \in \mathbb{R}^2$, define the Weyl operator
\begin{equation}\label{eq:weyl-op}
    W(a, b) = \exp\!\left(\frac{i}{\hbar_{\mathrm{econ}}} (a\hat{X} + b\hat{P})\right).
\end{equation}
Here the first parameter $a$ multiplies $\hat{X}$ (log-price) and the second
$b$ multiplies $\hat{P}$ (order-flow momentum). Under this convention and
$[\hat{X},\hat{P}]=i\hbar_{\mathrm{econ}}\mathbf{1}$, the Baker--Campbell--Hausdorff
formula yields the commutator
$[a\hat{X}+b\hat{P},\,a'\hat{X}+b'\hat{P}] = i\hbar_{\mathrm{econ}}(ab'-a'b)\mathbf{1}$,
giving the Weyl relation below.
\end{definition}

The Weyl operators \eqref{eq:weyl-op} satisfy the Weyl form of the CCR:
\begin{equation}\label{eq:weyl-ccr}
    W(a, b) W(a', b') = \exp\!\left(+\frac{i}{2\hbar_{\mathrm{econ}}} (ab' - a'b)\right)
    W(a + a', b + b').
\end{equation}

The Weyl operators generate the Weyl--Heisenberg algebra $\mathcal{A}$, which is the
$C^*$-algebra of market observables. The algebraic structure captures the
non-commutative nature of price and order flow: two Weyl operators commute if and
only if the corresponding phase-space area $ab' - a'b$ is an integer multiple of
$2\pi\hbar_{\mathrm{econ}}$.

\begin{definition}[Market observable algebra]\label{def:market-algebra}
The \emph{market observable algebra} $\mathcal{A}$ is the $C^*$-algebra generated by
the Weyl operators $\{W(a, b) : (a, b) \in \mathbb{R}^2\}$, equipped with the
involution $W(a, b)^* = W(-a, -b)$ and the unique $C^*$-norm determined by the CCR
\eqref{eq:weyl-ccr}.
\end{definition}

\subsection{The GNS Construction for Market States}

A market state is a positive, normalized, linear functional $\omega$ on
$\mathcal{A}$ that assigns to each observable its expected value. The
Gelfand--Naimark--Segal (GNS) construction represents this state as a vector in a
Hilbert space.

\begin{proposition}[GNS construction for market states {\normalfont\citep{takesaki2003theory}}]\label{def:gns}
Let $\omega$ be a faithful normal state on the market observable algebra
$\mathcal{A}$. Then there exists a triple
$(\mathcal{H}_\omega, \pi_\omega, |\Omega_\omega\rangle)$ consisting of:
\begin{enumerate}
    \item[(i)] a Hilbert space $\mathcal{H}_\omega$;
    \item[(ii)] a $*$-representation
        $\pi_\omega : \mathcal{A} \to \mathcal{B}(\mathcal{H}_\omega)$;
    \item[(iii)] a cyclic vector $|\Omega_\omega\rangle \in \mathcal{H}_\omega$ such
        that $\overline{\pi_\omega(\mathcal{A}) |\Omega_\omega\rangle} = \mathcal{H}_\omega$
        and $\omega(A) = \langle \Omega_\omega | \pi_\omega(A) | \Omega_\omega \rangle$
        for all $A \in \mathcal{A}$.
\end{enumerate}
This triple is unique up to unitary equivalence. The GNS construction is a standard
result of $C^*$-algebra theory; the inner product on $\mathcal{H}_\omega$ is defined by
$\langle \pi_\omega(A) \Omega_\omega \,|\, \pi_\omega(B) \Omega_\omega \rangle
= \omega(A^* B)$, and the representation is extended by continuity.
\end{proposition}

The GNS Hilbert space $\mathcal{H}_\omega$ is the space of all financial positions
that can be expressed as limits of portfolios of market observables. The cyclic
vector $|\Omega_\omega\rangle$ represents the \emph{vacuum state} of the market
---the ground state of the financial system in the absence of priced risk factors.

\subsection{Hedging Portfolios as Quantum Observables}

We now formulate the hedging problem within the GNS framework. Let $G(S_T)$ be the
payoff of a European option maturing at $T$, and let $Y = e^{-rT} G(S_T)$ be its
discounted value. Via the GNS construction, the discounted payoff is represented as
a vector $|Y\rangle \in \mathcal{H}_\omega$ defined by
\begin{equation*}
    |Y\rangle = \pi_\omega(e^{-rT} \hat{G}) |\Omega_\omega\rangle,
\end{equation*}
where $\hat{G} \in \mathcal{A}$ is the observable corresponding to the payoff
function.

\begin{definition}[Hedgeable subspace]\label{def:hedgeable-subspace}
Let $\mathcal{K} \subset \mathcal{H}_\omega$ be the closed linear subspace spanned
by the discounted asset prices available for hedging:
\begin{equation*}
    \mathcal{K} = \overline{\operatorname{span}}
        \{|S_{t_i}\rangle : t_i \in [0, T]\},
\end{equation*}
where $|S_{t_i}\rangle$ is the GNS vector corresponding to the discounted asset
price at time $t_i$. We call $\mathcal{K}$ the \emph{hedgeable subspace}.
\end{definition}

A self-financing hedging portfolio is represented by an operator $H \in \mathcal{A}$
whose GNS image $\pi_\omega(H)$ acts on $\mathcal{H}_\omega$. The optimal hedge is
the orthogonal projection of the payoff vector onto the hedgeable subspace.

\begin{theorem}[Quantum pricing formula]\label{thm:optimal-hedge}
Let $G(S_T)$ be the payoff of a European option maturing at $T$, let $r \geq 0$ be
the risk-free rate, and let $|\Omega_\omega\rangle \in \mathcal{H}_\omega$ be the
cyclic vector of the GNS representation. Then the arbitrage-free price of the option
at time $t \leq T$ is given by the quantum expectation value
\begin{equation}\label{eq:pricing-formula-basic}
    P_t = \langle \Omega_\omega | \pi_\omega(H) | \Omega_\omega \rangle
    \, e^{-r(T-t)},
\end{equation}
where $H \in \mathcal{A}$ is the self-adjoint operator representing the optimal
self-financing hedging portfolio, and $\pi_\omega(H)$ is its GNS representation.
Equivalently, in the State--Hedge representation
(Corollary~\ref{cor:state-hedge-pricing}), the optimal hedge representative
$h^* = \Pi_{\mathcal{K}_t} Y_t$ is the orthogonal projection of the payoff-state
vector $Y_t = g \Phi_t$ onto the hedgeable subspace $\mathcal{K}_t$, and the price
is the discounted real inner product
$P_t = \langle h^* \,|\, \Phi_t \rangle e^{-r(T-t)}$ on the real Hilbert space
$\mathcal{H}_\omega \simeq L^2(\mathbb{R}, \mathbb{R})$.
\end{theorem}

\begin{proof}
The proof proceeds in three steps.

\paragraph{Step 1: GNS representation of the payoff.}
The discounted payoff $Y = e^{-rT} G(S_T)$ is represented as a vector
$|Y\rangle \in \mathcal{H}_\omega$ via
\begin{equation*}
    |Y\rangle = \pi_\omega(e^{-rT} \hat{G}) |\Omega_\omega\rangle.
\end{equation*}

\paragraph{Step 2: Optimal hedge as orthogonal projection.}
The optimal hedge $|h^*\rangle$ is the unique vector in $\mathcal{K}$ that minimizes
the squared hedging error $\|Y - h\|_\omega^2$. By the Hilbert space projection
theorem, this is the orthogonal projection of $|Y\rangle$ onto $\mathcal{K}$:
\begin{equation*}
    |h^*\rangle = \Pi_{\mathcal{K}} |Y\rangle,
\end{equation*}
where $\Pi_{\mathcal{K}} \in \mathcal{B}(\mathcal{H}_\omega)$ is the orthogonal
projection operator. Since $|\Omega_\omega\rangle$ is cyclic for $\pi_\omega(\mathcal{A})$,
the orbit $\pi_\omega(\mathcal{A})|\Omega_\omega\rangle$ is dense in
$\mathcal{H}_\omega$. As $\mathcal{K}$ is a closed subspace and the projection
$\Pi_{\mathcal{K}}$ is continuous, there exists a unique self-adjoint operator
$H \in \mathcal{A}$ (in the closure of the GNS image) such that
$\pi_\omega(H) |\Omega_\omega\rangle = |h^*\rangle$.

\paragraph{Step 3: Discounting and no-arbitrage.}
The time-$t$ value of the hedge is the discounted expectation:
\begin{equation*}
    P_t = \langle \Omega_\omega | \pi_\omega(H) | \Omega_\omega \rangle e^{-r(T-t)}.
\end{equation*}
The market pricing functional
$\omega_t(A) = \langle \Omega_\omega | \pi_\omega(A) | \Omega_\omega \rangle$ is
positive and normalised by the GNS construction. Its faithfulness---and hence the
equivalence with the absence of arbitrage---is established in
Theorem~\ref{thm:fundamental} (Section~\ref{sec:no-arbitrage}), proved
independently of this formula.
\end{proof}

\begin{remark}[Relation to the classical BSM formula]\label{rem:bsm-limit}
In the classical limit $\hbar_{\mathrm{econ}} \to 0$, the canonical commutation
relation \eqref{eq:weyl-ccr} becomes $[\hat{X}, \hat{P}] = 0$, the observable
algebra $\mathcal{A}$ is commutative, and the GNS Hilbert space reduces to
$L^2(\mathbb{P})$. The hedgeable subspace $\mathcal{K}$ becomes the full space, so
$\Pi_{\mathcal{K}} = \mathbf{1}$ and the pricing formula collapses to
\begin{equation*}
    P_t^{\mathrm{BSM}} = e^{-r(T-t)} \mathbb{E}^{\mathbb{Q}}[G(S_T) \,|\, \mathcal{F}_t],
\end{equation*}
where $\mathbb{Q}$ is the unique risk-neutral measure. The quantum correction is
therefore a genuine generalization that contains BSM as the commutative special
case.
\end{remark}

\subsection{The Market Hamiltonian and Lindblad Dynamics}

The dynamics of market observables are generated by a self-adjoint Hamiltonian
operator $\hat{H}$. In the Heisenberg picture, the time evolution of an observable
$A \in \mathcal{A}$ is given by
\begin{equation*}
    \frac{d}{dt} A(t) = \frac{i}{\hbar_{\mathrm{econ}}} [\hat{H}, A(t)]
    + \mathcal{L}[A(t)],
\end{equation*}
where $\mathcal{L}$ is a Lindbladian superoperator that accounts for non-unitary,
dissipative dynamics \citep{lindblad1976generators,gorini1976completely}. The
Lindbladian takes the standard form
\begin{equation*}
    \mathcal{L}[A] = \sum_k \gamma_k
    \left( L_k^\dagger A L_k - \frac{1}{2} \{L_k^\dagger L_k, A\} \right),
\end{equation*}
where $\{L_k\}$ are Lindblad operators representing the sources of market friction
and $\gamma_k \geq 0$ are relaxation rates.

For the purpose of option pricing, we consider the Hamiltonian
\begin{equation}\label{eq:hamiltonian}
    \hat{H} = \frac{\hat{P}^2}{2M} + V(\hat{X}),
\end{equation}
where $M > 0$ is the \emph{market depth} parameter and $V(\hat{X})$ is a potential
function representing the confining effect of arbitrage bounds. The kinetic term
$\hat{P}^2/(2M)$ captures the cost of order flow: larger order flow generates larger
price impact, with the market depth $M$ determining the scale of this effect. The
potential $V(\hat{X})$ ensures that the log-price remains within economically
meaningful bounds.

The state of the market at time $t$ is described by a density operator $\rho_t$
acting on $\mathcal{H}_\omega$. In the Schr\"odinger picture, the evolution of the
state is governed by the Lindblad master equation:
\begin{equation}\label{eq:lindblad}
    \frac{d}{dt} \rho_t = -\frac{i}{\hbar_{\mathrm{econ}}} [\hat{H}, \rho_t]
    + \sum_k \gamma_k
    \left( L_k \rho_t L_k^\dagger - \frac{1}{2} \{L_k^\dagger L_k, \rho_t\} \right).
\end{equation}
The market risk state $|\Phi_t\rangle$ that appears in the pricing formula is
obtained from the density operator via a purification: there exists a Hilbert space
$\mathcal{H}_\omega \otimes \mathcal{H}_{\mathrm{anc}}$ such that
\begin{equation*}
    \rho_t = \operatorname{Tr}_{\mathrm{anc}}\!\left(|\Phi_t\rangle \langle \Phi_t|\right),
\end{equation*}
where $\operatorname{Tr}_{\mathrm{anc}}$ denotes the partial trace over an ancillary
space representing unobserved market degrees of freedom. The purified state
$|\Phi_t\rangle$ evolves under a stochastic Schr\"odinger equation driven by the
Lindblad dynamics, and the pricing formula \eqref{eq:pricing-formula-basic}
generalizes to
\begin{equation*}
    P_t = \langle \Phi_t | \pi_\omega(H) \otimes \mathbf{1}_{\mathrm{anc}} | \Phi_t \rangle
    \, e^{-r(T-t)},
\end{equation*}
which is the duality pairing that forms the computational basis for the empirical
methodology developed in Section~\ref{sec:empirical-calibration}.

% -----------------------------------------------------------------------
\subsection{Uniqueness of the Quantum Representation (Stone--von Neumann)}
\label{sec:stone-vn-uniqueness}
% -----------------------------------------------------------------------

A natural question is whether the adoption of the quantum Hilbert-space formalism
is a genuine consequence of the CCR, or merely one of many possible mathematical
representations. The Stone--von Neumann theorem resolves this: it shows that the
quantum mechanical Hilbert-space framework is the \emph{unique} irreducible
representation of the Weyl CCR algebra up to unitary equivalence.

\begin{proposition}[Uniqueness of the Weyl representation]
\label{prop:stone-vn}
Let $\mathcal{A}$ be the Weyl $C^*$-algebra generated by
$\{W(a,b): (a,b)\in\mathbb{R}^2\}$ satisfying \eqref{eq:weyl-ccr} with
$\hbar_{\mathrm{econ}} > 0$. Then every strongly continuous irreducible
unitary representation $(\mathcal{H}, W)$ of $\mathcal{A}$ is unitarily
equivalent to the Schr\"{o}dinger representation on $L^2(\mathbb{R},\mathbb{C})$:
\begin{equation*}
    (W(a,b)f)(x) = e^{i(ax + ab\hbar_{\mathrm{econ}}/2)/\hbar_{\mathrm{econ}}}
    f\!\left(x + b\right), \quad f \in L^2(\mathbb{R},\mathbb{C}).
\end{equation*}
In particular, the Hilbert-space structure $\mathcal{H}_\omega \simeq
L^2(\mathbb{R},\mathbb{C})$ is not an arbitrary modelling choice but the
\emph{unique irreducible carrier space} of the CCR derived in
Section~\ref{sec:ccr-derivation}.
\end{proposition}

\begin{proof}
This is the classical Stone--von Neumann theorem
\citep{vonNeumann1931,reed1972methods}. The key hypotheses---strong continuity
of $a \mapsto W(a,0)$ and $b \mapsto W(0,b)$ (Weyl regularity) and
irreducibility---are satisfied in the financial setting: strong continuity
follows from the continuity of the price-impact rule \eqref{eq:kyle-price-impact}
in the order size $q$, and irreducibility reflects the fact that a single
informed trader interacting with the full price-impact mechanism generates
the entire orbit of the observable algebra. The Schr\"{o}dinger representation
is therefore the canonical and unique framework for quantising the market
observables derived from Kyle's model.
\end{proof}

\begin{remark}[Implications for modelling]
Proposition~\ref{prop:stone-vn} addresses the concern that the quantum
probability framework is merely a formal rewriting of classical operator
algebra. Any other representation of $\mathcal{A}$ with $\hbar_{\mathrm{econ}}>0$
decomposes as a direct sum (or direct integral) of copies of the Schr\"{o}dinger
representation; the irreducible component is always the quantum Hilbert space.
The commutative ($\hbar_{\mathrm{econ}} = 0$) limit is excluded by
Proposition~\ref{prop:stone-vn}, as the Stone--von Neumann theorem does not apply
when the CCR degenerates to $[\hat{X},\hat{P}]=0$---consistent with the fact
that the frictionless market admits a classical ($L^2(\Omega,P)$) rather than
quantum representation.
\end{remark}
